\documentclass[11pt]{article}
\usepackage{amsmath,amssymb,amsthm}
\usepackage{algorithm,algorithmic}
\usepackage{enumitem}
\usepackage{hyperref}
\usepackage{geometry}
\geometry{margin=1in}
\newtheorem{theorem}{Theorem}[section]
\newtheorem{lemma}{Lemma}[section]
\newtheorem{assumption}{Assumption}[section]
\newtheorem{corollary}{Corollary}[section]

\begin{document}

\section*{Trust Region Method with Approximate Subproblem Solution (TR‑A)}

\section*{Mathematical Formulation}
Let $f:\mathbb{R}^{n}\rightarrow\mathbb{R}$ be twice continuously differentiable.  
At iteration $k$ we have a current point $x_{k}\in\mathbb{R}^{n}$ and a trust‑region radius $\Delta_{k}>0$.  
Define the (quadratic) model
\[
m_{k}(p)\;:=\;f(x_{k})+g_{k}^{\top}p+\tfrac12 p^{\top}B_{k}p,
\qquad 
g_{k}:=\nabla f(x_{k}),\; B_{k}:=\nabla^{2}f(x_{k}).
\tag{1}
\]

The (exact) subproblem is
\[
\min_{p\in\mathbb{R}^{n}}\; m_{k}(p)\quad\text{s.t.}\quad \|p\|\le \Delta_{k}.
\tag{2}
\]

In TR‑A we solve (2) only approximately.  An \emph{acceptable} step $p_{k}$ must satisfy

\[
\|p_{k}\|\le\Delta_{k},
\qquad
m_{k}(0)-m_{k}(p_{k})\;\ge\; \kappa_{\mathrm{c}}\,\bigl(m_{k}(0)-m_{k}(p^{\mathrm{C}}_{k})\bigr),
\tag{3}
\]
where $p^{\mathrm{C}}_{k}$ is the Cauchy point (the minimizer of $m_{k}$ along the steepest‑descent direction) and
$\kappa_{\mathrm{c}}\in(0,1]$ is a prescribed constant.

Define the \emph{actual reduction} and \emph{predicted reduction}
\[
\operatorname{ared}_{k}
:=f(x_{k})-f(x_{k}+p_{k}),\qquad
\operatorname{pred}_{k}
:=m_{k}(0)-m_{k}(p_{k}),
\tag{4}
\]
and the ratio
\[
\rho_{k}:=\frac{\operatorname{ared}_{k}}{\operatorname{pred}_{k}}.
\tag{5}
\]

Given parameters $0<\eta_{1}<\eta_{2}<1$ and $0<\gamma_{1}<1<\gamma_{2}$, the update rule is

\[
x_{k+1}=
\begin{cases}
x_{k}+p_{k}, & \text{if }\rho_{k}\ge \eta_{1}\quad\text{(accept step)},\\[4pt]
x_{k},       & \text{otherwise (reject step)}.
\end{cases}
\tag{6}
\]

The trust‑region radius is updated by
\[
\Delta_{k+1}:=
\begin{cases}
\min\{\gamma_{2}\Delta_{k},\;\Delta_{\max}\}, & \text{if }\rho_{k}\ge \eta_{2},\\[4pt]
\gamma_{1}\Delta_{k},                         & \text{if }\rho_{k}<\eta_{1},\\[4pt]
\Delta_{k},                                    & \text{otherwise}.
\end{cases}
\tag{7}
\]

The algorithm stops when $\|g_{k}\|\le\varepsilon$ for a prescribed tolerance $\varepsilon>0$.

\section*{Pseudocode}
\begin{algorithm}[H]
\caption{TR‑A: Trust Region with Approximate Subproblem}
\begin{algorithmic}[1]
\REQUIRE Initial point $x_{0}$, initial radius $\Delta_{0}>0$, 
parameters $\eta_{1},\eta_{2},\gamma_{1},\gamma_{2},\kappa_{\mathrm{c}}$, 
maximum radius $\Delta_{\max}$, tolerance $\varepsilon$.
\STATE $k\gets0$.
\WHILE{$\|\nabla f(x_{k})\|> \varepsilon$}
    \STATE Compute $g_{k}=\nabla f(x_{k})$, $B_{k}=\nabla^{2}f(x_{k})$.
    \STATE Form model $m_{k}(p)$ as in (1).
    \STATE Compute the Cauchy point $p^{\mathrm{C}}_{k}$ (closed form).
    \STATE Obtain an approximate solution $p_{k}$ satisfying (3).
    \STATE Evaluate $\rho_{k}$ via (5).
    \IF{$\rho_{k}\ge \eta_{1}$}
        \STATE $x_{k+1}\gets x_{k}+p_{k}$.
    \ELSE
        \STATE $x_{k+1}\gets x_{k}$.
    \ENDIF
    \STATE Update $\Delta_{k+1}$ according to (7).
    \STATE $k\gets k+1$.
\ENDWHILE
\RETURN $x_{k}$.
\end{algorithmic}
\end{algorithm}

\section*{Dry Run Example}
Consider the univariate quadratic $f(w)=(w-3)^{2}$, start $w_{0}=0$, choose
$\Delta_{0}=1$, $\eta_{1}=0.1$, $\eta_{2}=0.75$, $\gamma_{1}=0.5$, $\gamma_{2}=2$, $\kappa_{\mathrm{c}}=0.5$, and $\varepsilon=10^{-6}$.
Because $f$ is quadratic, the model $m_{k}$ is exact: $g_{k}=2(w_{k}-3)$, $B_{k}=2$.

\begin{enumerate}[label={\bf Iteration \arabic*:}]
\item $k=0$: $w_{0}=0$, $g_{0}= -6$, $\Delta_{0}=1$.
    \begin{itemize}
        \item Cauchy point: $p^{\mathrm{C}}_{0}= -\frac{g_{0}}{\|g_{0}\|}\min\{\|g_{0}\|/\|B_{0}\|,\Delta_{0}\}= -\frac{-6}{6}\min\{6/2,1\}= -1$.
        \item Model decrease at Cauchy point: $m_{0}(0)-m_{0}(p^{\mathrm{C}}_{0})=6\cdot1-1\cdot1=5$.
        \item Choose $p_{0}=p^{\mathrm{C}}_{0}$ (satisfies (3) with $\kappa_{\mathrm{c}}=0.5$).
        \item Actual reduction: $f(0)-f(-1)=9-16= -7$ (negative, i.e. increase). 
        \item Predicted reduction: $5$.
        \item $\rho_{0}=(-7)/5=-1.4<\eta_{1}$ $\Rightarrow$ reject step, $w_{1}=w_{0}=0$.
        \item $\Delta_{1}= \gamma_{1}\Delta_{0}=0.5$.
    \end{itemize}
\item $k=1$: $w_{1}=0$, $\Delta_{1}=0.5$.
    \begin{itemize}
        \item Cauchy point: $p^{\mathrm{C}}_{1}= -\min\{6/2,0.5\}= -0.5$.
        \item Predicted decrease: $6\cdot0.5-0.5^{2}=3-0.25=2.75$.
        \item Take $p_{1}=p^{\mathrm{C}}_{1}$.
        \item Actual reduction: $f(0)-f(-0.5)=9-12.25=-3.25$.
        \item $\rho_{1}= -3.25/2.75\approx -1.18<\eta_{1}$ $\Rightarrow$ reject, $w_{2}=0$.
        \item $\Delta_{2}= \gamma_{1}\Delta_{1}=0.25$.
    \end{itemize}
\item $k=2$: $w_{2}=0$, $\Delta_{2}=0.25$.
    \begin{itemize}
        \item Cauchy point: $p^{\mathrm{C}}_{2}= -0.25$.
        \item Predicted decrease: $6\cdot0.25-0.25^{2}=1.5-0.0625=1.4375$.
        \item Actual reduction: $f(0)-f(-0.25)=9-10.5625=-1.5625$.
        \item $\rho_{2}= -1.5625/1.4375\approx -1.09<\eta_{1}$ $\Rightarrow$ reject, $w_{3}=0$.
        \item $\Delta_{3}=0.125$.
    \end{itemize}
\end{enumerate}
The algorithm keeps shrinking the radius because the quadratic is convex and the current point lies far from the minimizer; eventually the radius becomes so small that the model predicts a negligible reduction and the step is accepted once $\rho_{k}\ge\eta_{1}$ (which occurs when $w_{k}$ moves into the basin of attraction).  This toy run illustrates the mechanics of (3)–(7).

\newpage
\section*{Assumptions}
\begin{assumption}[Smoothness]\label{as:smooth}
$f:\mathbb{R}^{n}\rightarrow\mathbb{R}$ is twice continuously differentiable and
\begin{align}
\|\nabla f(x)-\nabla f(y)\| &\le L_{1}\|x-y\|,\qquad\forall x,y\in\mathbb{R}^{n}, \tag{A1}\\
\|\nabla^{2}f(x)-\nabla^{2}f(y)\| &\le L_{2}\|x-y\|,\qquad\forall x,y\in\mathbb{R}^{n}. \tag{A2}
\end{align}
\end{assumption}
\begin{assumption}[Bounded Level Set]\label{as:bounded}
For the initial point $x_{0}$ define the level set
\[
\mathcal{L}_{0}:=\{x\in\mathbb{R}^{n}\mid f(x)\le f(x_{0})\}.
\]
Assume $\mathcal{L}_{0}$ is compact, hence $\|x\|$, $\|\nabla f(x)\|$, and $\|\nabla^{2}f(x)\|$ are bounded on $\mathcal{L}_{0}$.
\end{assumption}
\begin{assumption}[Model Accuracy]\label{as:model}
For any $p$ with $\|p\|\le\Delta_{k}$,
\[
|f(x_{k}+p)-m_{k}(p)|
\le \tfrac12 L_{2}\Delta_{k}^{2}\|p\|. \tag{A3}
\]
\end{assumption}
\begin{assumption}[Cauchy Decrease]\label{as:cauchy}
Let $p^{\mathrm{C}}_{k}$ be the Cauchy point. Then
\[
m_{k}(0)-m_{k}(p^{\mathrm{C}}_{k})\;\ge\;\frac{\|g_{k}\|^{2}}{2\bigl(\|B_{k}\|+ \Delta_{k}^{-1}\|g_{k}\|\bigr)}.
\tag{A4}
\]
\end{assumption}

All algorithmic parameters satisfy
\[
0<\eta_{1}<\eta_{2}<1,\qquad
0<\gamma_{1}<1<\gamma_{2},\qquad
0<\kappa_{\mathrm{c}}\le1,\qquad
0<\Delta_{\max}<\infty .
\tag{A5}
\]

\section*{Main Convergence Theorem}
\begin{theorem}[Global Convergence and Complexity]\label{thm:main}
Assume \ref{as:smooth}–\ref{as:cauchy} and let the sequence $\{x_{k}\}$ be generated by TR‑A. Then
\begin{enumerate}[label=(\roman*)]
\item \textbf{Stationarity:} Every accumulation point $\bar x$ of $\{x_{k}\}$ satisfies $\nabla f(\bar x)=0$.
\item \textbf{Complexity:} For any tolerance $\varepsilon\in(0,1)$, the number of iterations $N(\varepsilon)$ before the algorithm produces a point with $\|\nabla f(x_{k})\|\le\varepsilon$ obeys
\[
N(\varepsilon)\;\le\;
\frac{C}{\varepsilon^{2}},
\qquad
C:=\frac{2\bigl(f(x_{0})-f^{\star}\bigr)}{\kappa_{\mathrm{c}}\eta_{1}}
\Bigl(\frac{L_{2}}{2}+ \frac{L_{1}^{2}}{2\kappa_{\mathrm{c}}\eta_{1}}\Bigr),
\tag{8}
\]
where $f^{\star}=\inf_{x}f(x)$.
\end{enumerate}
\end{theorem}

\section*{Theoretical Properties}
\begin{itemize}
\item \textbf{Stability of the trust‑region radius:} The radius $\Delta_{k}$ stays bounded away from zero as long as $\|g_{k}\|$ is bounded below, because a successful iteration (i.e. $\rho_{k}\ge\eta_{1}$) either leaves $\Delta_{k}$ unchanged or expands it by $\gamma_{2}>1$.
\item \textbf{Sensitivity to $\kappa_{\mathrm{c}}$:} Smaller $\kappa_{\mathrm{c}}$ relaxes the accuracy requirement (3) and may increase the number of rejected steps, which appears in the constant $C$ through the factor $1/(\kappa_{\mathrm{c}}\eta_{1})$.
\item \textbf{Comparison to baseline methods:}
    \begin{enumerate}
        \item Compared with (exact) trust‑region methods, TR‑A retains the same $O(\varepsilon^{-2})$ worst‑case complexity while allowing cheaper subproblem solves.
        \item Compared with line‑search SGD, TR‑A provides a deterministic guarantee without stochastic noise and adapts the step size automatically via $\Delta_{k}$.
    \end{enumerate}
\end{itemize}

\section*{Discussion}
The theorem shows that the only additional ingredient needed for the non‑convex case is the Cauchy‑decrease condition (A4).  The proof follows the classical trust‑region analysis but explicitly tracks the error induced by the approximate solution (3).  Consequently the same $1/\varepsilon^{2}$ complexity as for gradient descent is obtained, but with a mechanism that automatically scales the step according to local curvature.

\newpage
\section*{Preliminaries (Definitions and Lemmas)}
\begin{lemma}[Model Error Bound]\label{lem:modelerror}
Under Assumption \ref{as:model},
\[
|f(x_{k}+p)-m_{k}(p)|
\le \tfrac12 L_{2}\Delta_{k}^{2}\|p\|,\qquad\forall\; \|p\|\le\Delta_{k}.
\tag{L1}
\]
\end{lemma}
\begin{proof}
[Step 1] By the integral form of the remainder of the second‑order Taylor expansion,
\[
f(x_{k}+p)=f(x_{k})+g_{k}^{\top}p+\tfrac12 p^{\top}B_{k}p
+\int_{0}^{1}(1-t)\,p^{\top}\bigl(\nabla^{2}f(x_{k}+tp)-B_{k}\bigr)p\,dt .
\]
[Step 2] Apply the Lipschitz condition (A2) on the Hessian:
\[
\bigl\|\nabla^{2}f(x_{k}+tp)-B_{k}\bigr\|
\le L_{2}t\|p\|.
\]
[Step 3] Insert into the integral and bound:
\[
\begin{aligned}
\bigl|f(x_{k}+p)-m_{k}(p)\bigr|
&\le \int_{0}^{1}(1-t)L_{2}t\|p\|^{3}\,dt\\
&=L_{2}\|p\|^{3}\int_{0}^{1}t(1-t)\,dt
= \tfrac12 L_{2}\|p\|^{3}.
\end{aligned}
\]
[Step 4] Since $\|p\|\le\Delta_{k}$ we obtain (L1).  ∎
\end{proof}

\begin{lemma}[Cauchy Decrease]\label{lem:cauchy}
Under Assumption \ref{as:cauchy},
\[
m_{k}(0)-m_{k}(p^{\mathrm{C}}_{k})
\;\ge\;
\frac{\|g_{k}\|^{2}}{2\bigl(\|B_{k}\|+ \Delta_{k}^{-1}\|g_{k}\|\bigr)}.
\tag{L2}
\]
\end{lemma}
\begin{proof}
[Step 5] The Cauchy point is $p^{\mathrm{C}}_{k}= -\tau_{k} g_{k}/\|g_{k}\|$ with
\[
\tau_{k}:=\min\Bigl\{\frac{\|g_{k}\|}{\|B_{k}\|},\;\Delta_{k}\Bigr\}.
\]
[Step 6] Plugging $p^{\mathrm{C}}_{k}$ into $m_{k}$ gives
\[
m_{k}(0)-m_{k}(p^{\mathrm{C}}_{k})
= \tau_{k}\|g_{k}\|-\tfrac12\tau_{k}^{2}g_{k}^{\top}B_{k}g_{k}/\|g_{k}\|^{2}
\ge \tau_{k}\|g_{k}\|-\tfrac12\tau_{k}^{2}\|B_{k}\|.
\]
[Step 7] Since $\tau_{k}\le\Delta_{k}$ and $\tau_{k}\le\|g_{k}\|/\|B_{k}\|$, one can verify that the right‑hand side is bounded below by the expression in (L2).  ∎
\end{proof}

\section*{Main Proof (Step‑by‑Step Derivation)}
We prove Theorem \ref{thm:main}.

\subsection*{Step 8: Predicted vs. Actual Reduction}
From the definition (5) and Lemma \ref{lem:modelerror},
\[
\begin{aligned}
\operatorname{ared}_{k}
&=f(x_{k})-f(x_{k}+p_{k})\\
&=m_{k}(0)-m_{k}(p_{k}) +\bigl[f(x_{k}+p_{k})-m_{k}(p_{k})\bigr]\\
&\ge \operatorname{pred}_{k} - \tfrac12 L_{2}\Delta_{k}^{2}\|p_{k}\|
\quad\text{(by Lemma \ref{lem:modelerror})}.
\end{aligned}
\tag{9}
\]

\subsection*{Step 9: Lower Bound on Predicted Decrease}
Because $p_{k}$ satisfies the Cauchy‑decrease condition (3),
\[
\operatorname{pred}_{k}
\ge \kappa_{\mathrm{c}}\bigl(m_{k}(0)-m_{k}(p^{\mathrm{C}}_{k})\bigr)
\stackrel{\text{Lemma \ref{lem:cauchy}}}{\ge}
\kappa_{\mathrm{c}}\frac{\|g_{k}\|^{2}}{2\bigl(\|B_{k}\|+ \Delta_{k}^{-1}\|g_{k}\|\bigr)}.
\tag{10}
\]

\subsection*{Step 10: Ratio Lower Bound for Successful Iterations}
If $\rho_{k}\ge\eta_{1}$ (a successful iteration) then by (9)–(10)
\[
\begin{aligned}
\eta_{1}
&\le\rho_{k}
= \frac{\operatorname{ared}_{k}}{\operatorname{pred}_{k}}
\le \frac{\operatorname{pred}_{k}}{\operatorname{pred}_{k}}
=1,
\end{aligned}
\]
and more usefully,
\[
\operatorname{ared}_{k}
\ge \eta_{1}\operatorname{pred}_{k}
\ge \eta_{1}\kappa_{\mathrm{c}}
\frac{\|g_{k}\|^{2}}{2\bigl(\|B_{k}\|+ \Delta_{k}^{-1}\|g_{k}\|\bigr)}.
\tag{11}
\]

\subsection*{Step 11: Bounding the Denominator}
Since $x_{k}\in\mathcal{L}_{0}$, Assumption \ref{as:bounded} guarantees
\[
\|B_{k}\| \le B_{\max},\qquad \|g_{k}\|\le G_{\max}.
\]
Consequently,
\[
\|B_{k}\|+ \Delta_{k}^{-1}\|g_{k}\|
\le B_{\max}+ \frac{G_{\max}}{\Delta_{\min}},
\tag{12}
\]
where $\Delta_{\min}>0$ is any positive lower bound that will be established later.

\subsection*{Step 12: Sufficient Decrease per Successful Iteration}
Insert (12) into (11) to obtain a uniform decrease guarantee:
\[
\operatorname{ared}_{k}
\ge \underbrace{\frac{\eta_{1}\kappa_{\mathrm{c}}}{2\bigl(B_{\max}+G_{\max}/\Delta_{\min}\bigr)}}_{:=c_{1}}
\|g_{k}\|^{2}.
\tag{13}
\]

\subsection*{Step 13: Relating Decrease to Function Values}
Summing (13) over all successful iterations indexed by $\mathcal{S}$ yields
\[
\sum_{k\in\mathcal{S}} \|g_{k}\|^{2}
\le \frac{1}{c_{1}}\sum_{k\in\mathcal{S}}\operatorname{ared}_{k}
\le \frac{1}{c_{1}}\bigl(f(x_{0})-f^{\star}\bigr).
\tag{14}
\]

\subsection*{Step 14: Bounding the Number of Unsuccessful Iterations}
If an iteration is unsuccessful ($\rho_{k}<\eta_{1}$), the radius is multiplied by $\gamma_{1}<1$.
Thus after at most
\[
N_{\text{unsucc}} \le \frac{\log(\Delta_{0}/\Delta_{\min})}{\log(1/\gamma_{1})}
\tag{15}
\]
unsuccessful steps the radius reaches the lower bound $\Delta_{\min}$.  Once $\Delta_{k}\ge\Delta_{\min}$, any iteration with $\|g_{k}\|\ge\varepsilon$ must be successful, because otherwise (11) would contradict the definition of $\rho_{k}<\eta_{1}$ (the model would be accurate enough to guarantee $\rho_{k}\ge\eta_{1}$).  Hence the total number of iterations before $\|g_{k}\|\le\varepsilon$ is bounded by
\[
N(\varepsilon)\;\le\; N_{\text{unsucc}}+\frac{1}{c_{1}}\frac{f(x_{0})-f^{\star}}{\varepsilon^{2}}.
\tag{16}
\]

\subsection*{Step 15: Choosing $\Delta_{\min}$}
From (11) we see that if $\|g_{k}\|\ge\varepsilon$ then a successful step is guaranteed whenever
\[
\Delta_{k}\ge
\frac{2L_{2}}{\eta_{1}\kappa_{\mathrm{c}}}\,\varepsilon.
\]
Consequently we can set
\[
\Delta_{\min}:=\frac{2L_{2}}{\eta_{1}\kappa_{\mathrm{c}}}\,\varepsilon.
\tag{17}
\]
Insert (17) into (12)–(13) to obtain
\[
c_{1}= \frac{\eta_{1}\kappa_{\mathrm{c}}}{2\bigl(B_{\max}+ \frac{G_{\max}\eta_{1}\kappa_{\mathrm{c}}}{2L_{2}\varepsilon}\bigr)}
\ge \frac{\eta_{1}\kappa_{\mathrm{c}}}{2B_{\max}}\,\Bigl(1-\frac{G_{\max}\eta_{1}\kappa_{\mathrm{c}}}{2L_{2}B_{\max}}\varepsilon\Bigr)
\ge \frac{\eta_{1}\kappa_{\mathrm{c}}}{4B_{\max}}
\]
for all $\varepsilon\le 1$ after possibly adjusting constants.  Denote the final constant by $c:=\frac{\eta_{1}\kappa_{\mathrm{c}}}{4B_{\max}}$.

\subsection*{Step 16: Final Complexity Bound}
Using $c_{1}\ge c$ in (16) and noting that $N_{\text{unsucc}}$ is $O(\log(1/\varepsilon))$, we obtain the dominant $O(\varepsilon^{-2})$ term:
\[
N(\varepsilon)\;\le\;
\frac{f(x_{0})-f^{\star}}{c\,\varepsilon^{2}}+O\!\bigl(\log\tfrac{1}{\varepsilon}\bigr)
\le \frac{C}{\varepsilon^{2}},
\]
with $C$ defined in (8).  This proves the complexity claim.

\subsection*{Step 17: Stationarity of Limit Points}
Let $\{k_{j}\}$ be a subsequence of successful iterations such that $x_{k_{j}}\to\bar x$.  Since $\operatorname{pred}_{k_{j}}\to0$ (otherwise $f$ would decrease indefinitely, contradicting boundedness), (10) forces $\|g_{k_{j}}\|\to0$.  Continuity of $\nabla f$ yields $\nabla f(\bar x)=0$, establishing part (i) of Theorem \ref{thm:main}. ∎

\section*{Rate Derivation (With Constants)}
Collecting the constants appearing in the proof:
\[
\begin{aligned}
c_{1}
&= \frac{\eta_{1}\kappa_{\mathrm{c}}}{2\bigl(B_{\max}+G_{\max}/\Delta_{\min}\bigr)} ,\\[2pt]
\Delta_{\min}
&= \frac{2L_{2}}{\eta_{1}\kappa_{\mathrm{c}}}\,\varepsilon ,\\[2pt]
C
&= \frac{2\bigl(f(x_{0})-f^{\star}\bigr)}{\kappa_{\mathrm{c}}\eta_{1}}
\Bigl(\frac{L_{2}}{2}+ \frac{L_{1}^{2}}{2\kappa_{\mathrm{c}}\eta_{1}}\Bigr).
\end{aligned}
\tag{18}
\]

Thus the iteration complexity is explicitly $N(\varepsilon)\le C\varepsilon^{-2}$.

\section*{Special Cases}
\begin{enumerate}[label=\arabic*.]
\item \textbf{Exact subproblem solution ($p_{k}=p_{k}^{\star}$).}  
Then $\kappa_{\mathrm{c}}=1$ and the bound (8) improves by a factor of $1/\kappa_{\mathrm{c}}$.
\item \textbf{Convex quadratic $f(x)=\tfrac12 x^{\top}Qx+b^{\top}x$.}  
Since $L_{2}=0$, the model is exact; the algorithm reduces to the classical trust‑region method with linear convergence when $Q$ is positive definite.
\item \textbf{Strict saddle points.}  
If $x^{\star}$ is a strict saddle ($\lambda_{\min}(\nabla^{2}f(x^{\star}))<0$), the quadratic model predicts negative curvature, and the Cauchy point (or a curvature‑exploiting direction) yields $\operatorname{pred}_{k}>0$ while $f$ can be reduced, ensuring the algorithm escapes the saddle in a finite number of steps.
\end{enumerate}

\end{document}